\subsection{RQ4: ablations}
\label{sec:rq_ablation}

We remove one component at a time under the grammar-kind view on BigVul (GGNN,
the same 10 folds as before) and report the paired change in AUC
(Table~\ref{tab:ablation}). Beyond the node views of Table~\ref{tab:direct},
which already ablate the contribution of text, anonymisation, kind and
operation type, two structural components are examined.

\paragraph{Edge types.} Removing the def-use edges changes AUC by
\numedgeNoDuDiff{} (\numedgeNoDuLo{} to \numedgeNoDuHi{}); removing the
sibling-order edges by \numedgeNoSibDiff{} (\numedgeNoSibLo{} to
\numedgeNoSibHi{}); keeping only the syntax-tree edges by
\numedgeAstOnlyDiff{} (\numedgeAstOnlyLo{} to \numedgeAstOnlyHi{}). We read
these as the contribution of the data-flow approximation and of source order on
top of the tree.

\paragraph{Operation-alphabet granularity.} Refining $\phi$ from eight to
sixteen operation types changes AUC by \numgranTOneSixDiff{}
(\numgranTOneSixLo{} to \numgranTOneSixHi{}) over \numgranTOneSixFolds{} matched
folds. The first submission reported a thirty-two-type run that rested on a
single fold; the claim that granularity does not matter was not supported by
that run and we do not repeat it. Both measurements we do have are consistent
with a null, as expected from Section~\ref{sec:rq_why}, since the classes that a
finer alphabet separates are rare: the memory-operation class fires on only
\numbvMemoryPct\% of nodes (Table~\ref{tab:node_dist}). The thirty-two-type
refinement was not re-run for this version because of its cost, so we make no
statement about it.

% Generated by experiments/emit_tables.py -- DO NOT EDIT BY HAND.
% Every value is read from experiments/results/*.json.

\begin{table}[t]
\centering
\caption{Ablations on BigVul (GGNN). First block: removing edge types from the
graph under the $\kappa$ view. Second block: refining the operation alphabet
under the $\kappa{+}\phi$ view, compared against the same view at
$|\mathcal{T}|=8$ on the same folds. $\Delta$AUC is paired, with 95\%
cluster-bootstrap interval and bootstrap $p$.}
\label{tab:ablation}
\resizebox{\linewidth}{!}{%
\begin{tabular}{lrrl}
\toprule
Variant & Folds & AUC & $\Delta$AUC vs reference [95\% CI] $p$ \\
\midrule
All edges (AST+SIB+DU), $\kappa$ & 10 & 0.691 & --- \\
\quad no def-use edges & 10 & 0.691 & $+0.000$ $[-0.012, +0.019]$ 0.758\,n.s. / 1.000 \\
\quad AST edges only & 10 & 0.699 & $+0.008$ $[-0.009, +0.027]$ 0.330\,n.s. / 0.375 \\
$\kappa{+}\phi$ with $|\mathcal{T}|=16$ vs $8$ & 10 & 0.691 & $+0.004$ $[-0.022, +0.019]$ 0.920\,n.s. / 0.625 \\
\bottomrule
\end{tabular}}
\end{table}

